\documentclass[a4paper,11pt,fleqn]{article}%fleqn pour aligner les equations à gauche
\usepackage{graphicx}
\usepackage[utf8]{inputenc}  
\usepackage[T1]{fontenc}
\usepackage{lscape}

\usepackage{tgheros,textcomp}%Police Arial
\renewcommand{\familydefault}{\sfdefault}

%\setlength{\mathindent}{0pt}%indentation dans math = 0

\usepackage[top = 1cm, bottom = 2cm, left = 1cm, right = 1cm]{geometry}
\usepackage{titlesec}
%\setcounter{secnumdepth}{4}
\usepackage{xcolor}
\usepackage{amssymb} %double flèche de réaction
\usepackage{mathrsfs}%farraday
\usepackage{xtab}%tableau sur plusieurs pages
\usepackage{esvect}%grands vecteurs

\usepackage[version=4]{mhchem} %equations chimiques

\usepackage{array}
\usepackage{chemfig}
\usepackage{multirow}
%\usepackage{sectsty}

\usepackage{caption}
\usepackage{adjustbox}
\usepackage[labelformat=simple]{subfig}
\usepackage{float}


\usepackage{amsmath}
\usepackage{amssymb}
\usepackage[cal=boondoxo]{mathalfa} 
\usepackage{enumitem}
\usepackage{lastpage}
\usepackage{tcolorbox}
\usepackage{siunitx}
\usepackage[european, straightvoltages, RPvoltages]{circuitikz}
\usepackage{tikz}
\usetikzlibrary{shapes.geometric}
\usetikzlibrary{calc}
\usetikzlibrary{automata,positioning}
\usepackage{multicol}
\usepackage{eqnarray}
\usepackage{tabularx,colortbl}%colorer fond cellules

\sisetup{
	detect-all,
	output-decimal-marker={,},
	group-minimum-digits = 3,
	group-separator={~},
	number-unit-separator={~},
	inter-unit-product={~}
} %setup pour SIunitx

\usepackage{listings}%pour insérer le code python
\definecolor{codegreen}{rgb}{0,0.6,0}
\definecolor{codegray}{rgb}{0.5,0.5,0.5}
\definecolor{codepurple}{rgb}{0.58,0,0.82}
\definecolor{backcolour}{rgb}{0.95,0.95,0.92}


%\graphicspath{ {/home/rweye/Documents/Enseignement/2023/Cours/TSPE/Eval/Ressources/} }

\titleformat{\section}
{\titlerule
\vspace{.8ex}%
\normalfont \bfseries}
{\thesection :}{.5em}{\bfseries}

\renewcommand{\thesection}{Exercice \arabic{section}}
\titlespacing{\section}{0pt}{1em}{1em}
\renewcommand{\thesubsection}{\hspace{1cm}\alph{subsection}.}
\renewcommand\thesubfigure{\textbf{Document \arabic{subfigure} :}}

\title{\vspace{-1cm}\large 
	\begin{tabular}{|m{5cm}|>{\centering}m{8cm}|>{\centering\arraybackslash}m{4cm}|}
		\hline
		Nom : & \textbf{\textsc{DS TSPE N°3}} &  Note : \\
		Prénom : & & $\_\_\_$ \textbf{20} \\ \hline
\end{tabular}}
\date{}

\newpagestyle{mystyle}{%
	\setfoot{ROSALIE}{\thepage\,$|$\,\pageref{LastPage}}{}
}
\renewpagestyle{plain}{%
	\setfoot{ROSALIE}{\thepage\,$|$\,\pageref{LastPage}}{}
}

\pagestyle{mystyle}

\newenvironment{itemize*}%
{\begin{itemize}[label=$-$]%
		\setlength{\itemsep}{0pt}%
		\setlength{\parskip}{0pt}}%
	{\end{itemize}}

\newenvironment{itemize**}%
{\begin{itemize}[label=$\_\_$,labelindent=0pt, leftmargin=*,topsep=0pt]%
		\setlength{\itemsep}{0pt}%
		\setlength{\parskip}{0pt}}%
	{\end{itemize}}

\renewcommand{\arraystretch}{1.5}


\begin{document}
\maketitle
\vspace{-2.5cm}

\noindent \begin{minipage}[t]{.9\linewidth}
\section{ Une réparation  ($\_\_\_$9)}
Un joueur de football doit tirer un coup franc, il se prépare à frapper et tire le ballon à une vitesse initiale de \SI{29,0}{\meter\per\second}. Le ballon part avec un angle par rapport au sol $\alpha $ = \SI{16}{\degree}. Le joueur a pour objectif de faire passer son tir au dessus du mur formé par les joueurs adverses situés à \SI{9,15}{m}. Le mur est estimé à \SI{1,95}{m} de haut. On négligera (à tord) l'effet des frottements pour l'ensemble de cet exercice.

\begin{enumerate}
	\item Réaliser un schéma de la situation initiale.
	\item Énoncer la deuxième loi de Newton et en déduire les équations correspondant à $a_x(t)$ et $a_z(t)$.
	\item Exprimer $x(t)$ et $z(t)$.
	\item Déterminer si le tir peut passer au dessus du mur.
\end{enumerate}

Peu après le tir, l'arbitre siffle pour indiquer la mi-temps. Les spécifications du sifflet indiquent que celui-ci produit un niveau d'intensité sonore de \SI{115}{dB} à une distance de \SI{15,0}{cm}.
\begin{enumerate}[resume]
	\item Déterminer l'intensité sonore associée au son du sifflet pour une distance d'écoute de \SI{15,0}{cm}.
	\item Déterminer la puissance associée au sifflet.
	\item Un joueur est situé de l'autre coté du terrain, à environ \SI{50,0}{m} de l'arbitre. Calculer le niveau d'intensité sonore perçu par ce joueur.
	\item Déterminer l'atténuation correspondante.
\end{enumerate}

\textbf{Données :}
\begin{itemize}
	\item Masse du ballon  : \SI{450}{g}
	\item Valeur de la pesanteur terrestre : $g$ = \SI{9.81}{\newton\per\kilo\gram}
\end{itemize}
\end{minipage}
\begin{minipage}[t]{.1\linewidth}
\vspace{8em} 
\begin{itemize}[label=$\_\_\_$]%
	\item 1
	\item 1
	\item 1
	\item 2
\end{itemize}
\vspace{2em} 
\begin{itemize}[label=$\_\_\_$]%
	\item 1
	\item 1
	\item 1\\
	\item 1
\end{itemize}
\end{minipage}

\noindent \begin{minipage}[t]{.9\linewidth}
\section{Lancement d'un avion de chasse ($\_\_\_$10)}
Sur un porte avion, la distance normalement nécessaire pour un avion de chasse ne peut pas être atteinte car le bâtiment aurait des dimensions trop importantes. Pour permettre à des avions de décoller sur une plus petite distance, les porte-avions disposent d'une catapulte aéronautique. Il s'agit d'un système permettant à un avion d'accélérer suffisamment pour atteindre leur vitesse de décollage.\\

\begin{enumerate}
	\item On considérera deux positions de l'avion : le point A correspond à sa position de départ, le point B correspond au point final de son « roulage » sur la piste (les roues touchent encore la poste). Représenter, sur un schéma la situation avec un système d'axes que vous préciserez les deux points A et B et la distance $d$ séparant ces deux points. Lors d’un décollage un pilote est soumis à une accélération correspondant à un peu plus de $3g$.
	\item Décrire le mouvement dans le référentiel du porte-avion.
	\item A quelle condition peut-on dire qu'il s'agit d'un référentiel galiléen.
	\item En se plaçant au point B de la trajectoire, établir un bilan des forces s'appliquant sur l'avion.
	\item Établir l'expression de l'accélération $a_x(t)$.
\end{enumerate}

On souhaite déterminer si la valeur de l'accélération subie par le pilote est bien celle proposée par l'énoncé.

\begin{enumerate}[resume]
	\item Calculer la variation d'énergie cinétique entre les points A et B.
	\item En déduire, grâce au théorème de l'énergie cinétique que la force de traction de la catapulte vaut $F$ = {643}{kN}
	\item Calculer la valeur de l'accélération $a_x(t)$ et vérifier l'information donnée dans l'énoncé.
\end{enumerate}

Pour déterminer la vitesse de l'avion au point B, on utilise un radar (émetteur et récepteur) qui émet une onde ultra-sonore d'une fréquence de \SI{500}{kHz} que l'on place très proche du point B.
\begin{enumerate}[resume]
	\item Décrire le phénomène observé lors d'un décollage.
	\item Montrer que le décalage de fréquence obtenu est de \SI{204}{kHz}
\end{enumerate}
\textbf{Données : }
\begin{itemize}
	\item Masse d’un avion de chasse chargé : \SI{20.0}{t}
	\item Longueur de la piste de décollage : \SI{75.0}{m}
	\item Vitesse de décollage : \SI{250}{\kilo\meter\per\hour}
	\item Vitesse du son dans l'air : \SI{340}{\meter\per\second}
	\item Décalage de fréquence dans notre cas : $|\Delta f| = \dfrac{2\cdot f_{em} \cdot v}{v_{son}}$
\end{itemize}
\textbf{Respect des chiffres significatifs : $\_\_\_$ 1pt}
\end{minipage}
\begin{minipage}[t]{.1\linewidth}
\vspace{10.5em}
\begin{itemize}[label=$\_\_\_$]%
	\item 1 \vspace{4em}
	\item 0.5
	\item 1
	\item 1
	\item 1
\end{itemize}
\vspace{2em} 
\begin{itemize}[label=$\_\_\_$]%
	\item 1
	\item 1.5\\
	\item 1
\end{itemize}
\vspace{2em} 
\begin{itemize}[label=$\_\_\_$]%
	\item 1
	\item 1
\end{itemize}
\end{minipage}

\end{document}